\chapter*{LỜI CẢM ƠN}
\addcontentsline{toc}{chapter}{LỜI CẢM ƠN}

Em xin bày tỏ lòng biết ơn sâu sắc tới GS. TS. Nguyễn Linh Trung và Thượng tá, TS. Nguyễn Thành Trung, Phó Chủ nhiệm Khoa Trang bị, Bệnh viện 108. Sự định hướng tận tình, những góp ý quý báu và các trao đổi nghiêm túc của các thầy đã giúp em từng bước xác định hướng đi, hiểu rõ hơn yêu cầu khoa học của đề tài và có thêm niềm tin để kiên trì hoàn thành khóa luận.

Em xin chân thành cảm ơn Cử nhân Nguyễn Phương Trang đã luôn sẵn lòng hỗ trợ, hướng dẫn và trao đổi với em trong suốt quá trình thực hiện đề tài. Những góp ý cụ thể cùng sự đồng hành ấy đã giúp em tháo gỡ nhiều vướng mắc và tiếp thêm động lực mỗi khi công việc nghiên cứu trở nên khó khăn.

Em xin trân trọng cảm ơn ông Sepehr Aghajanian, tác giả của bài báo nền tảng được tham khảo trong khóa luận, đã dành thời gian hồi đáp email và giải đáp những câu hỏi của em. Sự cởi mở và tận tâm ấy giúp em hiểu sâu hơn công trình nghiên cứu mà mình kế thừa, đồng thời khiến em thêm trân trọng tinh thần chia sẻ trong cộng đồng khoa học.

Em biết ơn tập thể Lab Avitech đã tạo điều kiện sử dụng tài nguyên máy chủ phục vụ việc huấn luyện mô hình và thực hiện các thí nghiệm. Sự hỗ trợ thiết thực về hạ tầng tính toán đã giúp em có thể theo đuổi quá trình thực nghiệm và hoàn thành nghiên cứu.

Em xin cảm ơn các thầy cô Trường Đại học Công nghệ, Đại học Quốc gia Hà Nội đã trao cho em nền tảng kiến thức quý báu. Em cũng xin ghi nhận đóng góp của các nhà nghiên cứu, cơ sở lâm sàng, nhà tài trợ và người tham gia Alzheimer's Disease Neuroimaging Initiative (ADNI), những người đã cùng xây dựng và duy trì nguồn dữ liệu có giá trị cho nghiên cứu này.

Sau cùng, em xin gửi lời biết ơn chân thành tới gia đình và những người thân yêu vì tình thương, niềm tin và sự động viên bền bỉ. Sự ủng hộ ấy là điểm tựa giúp em vững lòng trước những thử thách và đi đến những dòng cuối cùng của khóa luận.
